\documentclass[11pt]{article}
\usepackage[T1]{fontenc}
\usepackage[utf8]{inputenc}
\usepackage{lmodern}
\usepackage[a4paper,margin=28mm]{geometry}
\usepackage{amsmath,amssymb,amsthm,booktabs,microtype}
\usepackage[hidelinks]{hyperref}
\usepackage{url}
\hypersetup{
  pdftitle={Positivity of one-sided fractional Gaussian noise predictors for five increments},
  pdfauthor={Xiaojian Zeng},
  pdfsubject={A finite-case theorem with an exact rational certificate},
  pdfkeywords={fractional Gaussian noise, prediction, positivity, Bernstein certificate}}
\newtheorem{theorem}{Theorem}
\newtheorem{lemma}[theorem]{Lemma}
\newtheorem{proposition}[theorem]{Proposition}
\newcommand{\R}{\mathbb{R}}
\newcommand{\E}{\mathbb{E}}
\newcommand{\dd}{\,\mathrm{d}}
\DeclareMathOperator{\Var}{Var}
\DeclareMathOperator{\Cov}{Cov}
\title{Positivity of one-sided fractional Gaussian noise\protect\\ predictors for five increments}
\author{Xiaojian Zeng\\
  \small PBC School of Finance, Tsinghua University, 100190\\
  \small \href{mailto:xiaojianzeng@foxmail.com}{xiaojianzeng@foxmail.com}}
\date{}
\begin{document}
\maketitle

\begin{abstract}
For fractional Gaussian noise with Hurst index $1/2<H<1$, we prove that
every coefficient of the projection of the first increment onto the next
four increments is strictly positive. Together with the preceding cases,
this establishes one-sided predictor positivity for $2\le n\le5$ in the
usual $n$-increment notation. The original conjecture for arbitrary $n$
remains unresolved. The main step treats the penultimate coefficient at
$n=5$, singled out for numerical treatment in earlier work. Its Cramer
numerator is a thirty-term exponential sum. After removing its exact
endpoint zeros, a rational Taylor--Bernstein certificate proves a uniform
lower bound of $27/10$ for the normalized sum, including its analytic
endpoint continuations. All certificate data and a standard-library
Python verifier are supplied.
\end{abstract}

\section{The finite case and its context}

Put $t=2H\in(1,2)$ and write the fractional Gaussian noise covariance as
\begin{equation}\label{eq:rho}
 \rho_0=1,\qquad
 \rho_j=\frac{(j+1)^t-2j^t+|j-1|^t}{2}\quad(j\ge1).
\end{equation}
For $n\ge2$, define $\Gamma_n^k$, $2\le k\le n$, by
\begin{equation}\label{eq:normal}
 \sum_{k=2}^{n}\Gamma_n^k\rho_{|\ell-k|}=\rho_{\ell-1},
 \qquad 2\le\ell\le n.
\end{equation}
These are the coefficients of $\E(X_1\mid X_2,\ldots,X_n)$ for centered
Gaussian increments with covariance~\eqref{eq:rho}. The all-$n$ positivity
conjecture was posed by Mishura, Ralchenko and Schilling
\cite{Mishura2022} (Conjecture~(A) in the first arXiv version) and restated
as Conjecture~3 by Bodnarchuk, Mishura and Ralchenko
\cite{Bodnarchuk2025}. We prove the following restricted statement.

For $m\ge1$, let
\[
 T_m(t)=(\rho_{|i-j|})_{i,j=1}^{m},\qquad
 b_m(t)=(\rho_1,\ldots,\rho_m)^{\mathsf T},\qquad D_m(t)=\det T_m(t).
\]
Let $N_{nk}(t)$ be the determinant obtained by replacing column $k-1$
of $T_{n-1}(t)$ by $b_{n-1}(t)$. Thus $\Gamma_n^k=N_{nk}/D_{n-1}$.

\begin{theorem}\label{thm:main}
For every real $1/2<H<1$, every integer $2\le n\le5$, and every
$2\le k\le n$, one has $\Gamma_n^k(H)>0$.
Moreover, set $P(t)=16N_{54}(t)$, so that column~3 of $T_4$ is replaced
by $b_4$. The function
\begin{equation}\label{eq:margin}
 \frac{P(t)}{(t-1)(2-t)^3}
\end{equation}
extends analytically across both endpoints of $[1,2]$, and its continued
value is strictly greater than $27/10$ at every point of that interval.
\end{theorem}

Proposition~1 of \cite{Bodnarchuk2025} already states the $n=4$
inequality; its proof combines analytic subintervals with a graph.
Remark~4 of that paper identifies $\Gamma_5^4$ as a more difficult
penultimate coefficient and discusses it numerically. The contribution
here is a proof over the entire Hurst interval for this specified $n=5$
case, together with the quantitative bound~\eqref{eq:margin}. The
smaller cases and the general prediction machinery are supporting facts.
A bounded primary-source search through 13 September 2026 found no
equivalent full-interval $n=5$ certificate; this is not an exhaustive
priority claim. Cai's version~2 preprint \cite{Cai2026}, Theorems~2.2--2.3,
studies second-order asymptotics of the partial autocorrelation function.
Those asymptotic statements concern the last, or reflection, coefficient
and do not determine the finite interior coefficient treated here.

The proof uses elementary projection identities, an exponential zero
count, and a Taylor polynomial expressed in the Bernstein basis. The
last step is related to established Bernstein-form computation
\cite{FaroukiRajan1988}. Exact rational Bernstein controls also occur in
Zeng's restricted matrix-positivity certificates
\cite[Theorem~10 and equations~(8.25)--(8.27)]{Zeng2026Lens}.
For the present exponential sum, an explicit remainder estimate is
needed in addition to the polynomial controls. No general approximation
method is claimed as a contribution of this note.

\section{Covariance inequalities and the outer coefficients}

\begin{lemma}\label{lem:covariance}
All finite matrices $T_m(t)$ are strictly positive definite for $0<t<2$.
For $1<t<2$,
\[
 1=\rho_0>\rho_1>\rho_2>\cdots>0,
 \qquad \rho_j^2<\rho_{j-1}\rho_{j+1}\quad(j\ge1).
\]
In particular, $\rho_{j+1}/\rho_j$ increases strictly for $j\ge0$.
\end{lemma}

\begin{proof}
Let $C_t=\int_0^\infty(1-\cos u)u^{-1-t}\dd u$.
This integral is finite and positive for $0<t<2$. Scaling gives an
integral representation of $|j|^t$; taking its centered second difference
gives, also at $j=0$,
\[
 \rho_{|j|}=\frac1{C_t}\int_0^\infty
       \cos(ju)(1-\cos u)u^{-1-t}\dd u.
\]
Consequently, for a nonzero real vector $v$,
\[
 v^{\mathsf T}T_m v=\frac1{C_t}\int_0^\infty
   \left|\sum_{j=1}^m v_j e^{iju}\right|^2
   (1-\cos u)u^{-1-t}\dd u>0.
\]
Indeed, the weight is positive apart from a discrete set, and a nonzero
trigonometric polynomial cannot vanish on an open interval.

For $j\ge1$ and $1<t<2$, the integral form of a second difference yields
\begin{equation}\label{eq:second-difference}
 \rho_j=\frac{t(t-1)}2\int_{-1}^1
             (1-|u|)(j+u)^{t-2}\dd u.
\end{equation}
The singularity at $j=1,u=-1$ is integrable. This proves strict decrease
and positivity for $j\ge1$, while $\rho_1=2^{t-1}-1\in(0,1)$.
For $j\ge2$, strict log-convexity of $z^{t-2}$ gives the pointwise
geometric-mean inequality between its values at $j-1+u,j+u,j+1+u$.
Integration and Cauchy--Schwarz imply
$\rho_j^2<\rho_{j-1}\rho_{j+1}$. At $j=1$, direct substitution gives
\begin{equation}\label{eq:special-identity}
 \rho_2-\rho_1^2=\frac{\rho_1-\rho_3}{2}>0.
\end{equation}
Dividing these strict inequalities by positive factors proves the ratio
statement. These covariance properties are also developed
in~\cite{Mishura2022}.
\end{proof}

\begin{lemma}\label{lem:edges}
If $\Gamma_n^k>0$ for every $2\le k\le n$, then
$\Gamma_{n+1}^{2}>0$ and $\Gamma_{n+1}^{n+1}>0$.
\end{lemma}

\begin{proof}
Write $g_k=\Gamma_n^k$. Residualize $X_2$ against
$X_3,\ldots,X_{n+1}$. Stationarity identifies its projection coefficients
with the $g_k$. The coefficient of $X_2$ in the enlarged predictor has
a positive residual-variance denominator and numerator
\[
 \rho_1-\sum_{k=2}^n g_k\rho_k
   =\sum_{k=2}^n g_k(\rho_{k-2}-\rho_k)>0,
\]
where the equality uses~\eqref{eq:normal} at $\ell=2$.
For the other edge, residualize $X_{n+1}$ against $X_2,\ldots,X_n$.
Time reversal gives
\begin{equation}\label{eq:last}
 \Gamma_{n+1}^{n+1}=
 \frac{\rho_n-\sum_{k=2}^n g_k\rho_{n+1-k}}
      {1-\sum_{k=2}^n g_k\rho_{k-1}}.
\end{equation}
The denominator is again a strictly positive residual variance.
Since $\rho_{n-1}=\sum_{k=2}^n g_k\rho_{n-k}$, its numerator is
\[
 \sum_{k=2}^n g_k\rho_{n-k}
 \left(\frac{\rho_n}{\rho_{n-1}}
              -\frac{\rho_{n+1-k}}{\rho_{n-k}}\right)>0
\]
by Lemma~\ref{lem:covariance}. This is the conditional last-coefficient
argument of \cite[Proposition~2]{Bodnarchuk2025}. Positivity of the
entire preceding row is an explicit hypothesis in both edge steps.
\end{proof}

\section{The smaller interior signs}

Put
\begin{equation}\label{eq:qK}
 \begin{split}
 q&=\rho_2+\rho_2^2-\rho_1^2-\rho_1\rho_3,\\
 K&=(1-\rho_2)(1-\rho_4)-(\rho_1-\rho_3)^2.
 \end{split}
\end{equation}
Expanding the defining Cramer determinants gives
\begin{equation}\label{eq:factorizations}
 N_{43}=(1-\rho_2)q,\qquad N_{53}=qK.
\end{equation}
The covariance matrix of $X_1-X_5$ and $X_2-X_4$ is
\[
 2\begin{pmatrix}
 1-\rho_4&\rho_1-\rho_3\\
 \rho_1-\rho_3&1-\rho_2
 \end{pmatrix}.
\]
Its coefficient vectors are linearly independent, so
Lemma~\ref{lem:covariance} makes this matrix positive definite.
Its determinant is $4K$, hence $K>0$.

To prove $q>0$, substitution in~\eqref{eq:rho} gives
\begin{equation}\label{eq:F}
 F(t):=4q=9^t-8^t-2\cdot6^t+4\cdot4^t-2\cdot2^t-1.
\end{equation}
We use the following form of the generalized Descartes rule. Background
on zero counting by sign changes and Rolle's theorem is given
in~\cite{Jameson2006}; the needed argument is included here.

\begin{lemma}\label{lem:zeros}
If $\lambda_1<\cdots<\lambda_r$ are real and each $c_i\ne0$, the real
zeros of $\sum_{i=1}^r c_i e^{\lambda_i t}$, counted with multiplicity,
number at most the sign variations in $(c_1,\ldots,c_r)$.
\end{lemma}

\begin{proof}
Induct on the number $V$ of variations. For $V=0$, every summand has
the same sign. If $V>0$, choose $\lambda$ strictly between two consecutive
exponents at a sign change. Differentiating the exponential sum after
multiplication by $e^{-\lambda t}$ changes every coefficient to
$(\lambda_i-\lambda)c_i$. All signs to the left of $\lambda$ are reversed,
so exactly one variation disappears. The derivative has at most $V-1$
zeros by induction. Rolle's theorem, including multiplicities, bounds the
original number by $V$. Apply this on any finite interval containing a
proposed collection of zeros.
\end{proof}

In increasing order of bases, the coefficients of $F$ are
$-1,-2,4,-2,-1,1$, with three variations. Direct calculation gives
\[
 F(1)=F(2)=0,\quad F(3)=24,\quad
 F'(1)=2\log(27/16)>0,\quad F'(2)=18\log(243/256)<0.
\]
There is a third zero in $(2,3)$, because $F$ is negative just to the
right of $2$. These zeros exhaust the bound in Lemma~\ref{lem:zeros}.
It follows that $F>0$ on $(1,2)$, hence $q>0$. Thus
\eqref{eq:factorizations} proves $\Gamma_4^3>0$ and $\Gamma_5^3>0$.
The $n=4$ statement itself is already in
\cite[Proposition~1]{Bodnarchuk2025}.

\section{The penultimate cofactor at five increments}

The remaining sign is that of $P(t)=16N_{54}(t)$. Its exact expansion is
\begin{equation}\label{eq:P}
 P(t)=\sum_b c_b b^t,
\end{equation}
with the thirty integer pairs $(b,c_b)$ in Table~\ref{tab:terms}.
They can be reconstructed by expanding the defining determinant using
the polynomial entries
\[
 2\rho_0=2,\quad 2\rho_1=a-2,\quad
 2\rho_2=b-2a+1,\quad 2\rho_3=a^2-2b+a,\quad
 2\rho_4=c-2a^2+b,
\]
where in this identity only $a=2^t,b=3^t,c=5^t$.
The factor $16$ comes from scaling all four matrix columns by $2$.
Differentiation of~\eqref{eq:P} gives the exact identities
\begin{equation}\label{eq:endpoints}
 P(1)=0,\qquad P(2)=P'(2)=P''(2)=0,
 \qquad P'(1)=80\log2-48\log3.
\end{equation}
For example, each derivative at an endpoint is a polynomial in
$\log2,\log3,\log5$ with integer coefficients. Expanding these formal
polynomials verifies each zero coefficient exactly. This does not require
algebraic independence of the three logarithms. Since $P$ is entire,
the zeros in~\eqref{eq:endpoints} justify the analytic continuations
in Theorem~\ref{thm:main}.

\subsection{Removing the endpoint zeros before approximation}

Set $x=2-t$ and $A_k=(-1)^kP^{(k)}(2)/k!$. Then
\[
 P(2-x)=x^3G(x),\qquad
 G(x)=\sum_{k=3}^{\infty}A_kx^{k-3},\qquad G(1)=0.
\]
The continued expression~\eqref{eq:margin} is $R(x)=G(x)/(1-x)$.
For $N=26$, define
\begin{equation}\label{eq:RN}
 R_N(x)=\frac{G_N(x)-G_N(1)}{1-x}
       =-\sum_{k=4}^N A_k\sum_{j=0}^{k-4}x^j,
 \qquad G_N(x)=\sum_{k=3}^N A_kx^{k-3}.
\end{equation}
This is an exactly divisible polynomial of degree $22$; an approximate
endpoint zero is never divided by $1-x$. Absolute convergence and
$G(1)=0$ give, uniformly on $[0,1]$,
\begin{equation}\label{eq:tail}
 |R(x)-R_N(x)|\le\sum_{k=N+1}^{\infty}(k-3)|A_k|.
\end{equation}
This also controls the continued endpoint $x=1$.

Table~\ref{tab:terms} gives $b\le108$, $\log b<5$, and
$S:=\sum_b|c_b|b^2=212096$. Therefore $|A_k|\le S5^k/k!$.
For $k\ge27$, the ratio of consecutive majorant terms
$(k-3)5^k/k!$ is at most $5/24$. Summing a geometric majorant in
\eqref{eq:tail} yields
\begin{equation}\label{eq:B}
 |R-R_N|\le
 \frac{212096\cdot24\cdot5^{27}}{27!}\frac1{1-5/24}
 =\frac{790119171142578125}{179590728782941572096}
 <\frac1{200}.
\end{equation}

\subsection{Rational Bernstein controls}

All finite arithmetic bounds are rational. For $p\in\{2,3,5\}$ put
$z_p=(p-1)/(p+1)$ and
\begin{equation}\label{eq:logs}
 L_p=2\sum_{j=0}^{119}\frac{z_p^{2j+1}}{2j+1},\qquad
 0<\log p-L_p\le
       \frac{2z_p^{241}}{241(1-z_p^2)}.
\end{equation}
The estimate follows by bounding every omitted denominator below by
$241$ and summing a geometric series. Outward rounding at denominator
$10^{30}$ gives the log intervals in Table~\ref{tab:logs}.
Prime factorization then encloses every $\log b$ in~\eqref{eq:P}.
Substitute these positive intervals into
\begin{equation}\label{eq:A}
 A_k=\frac{(-1)^k}{k!}\sum_b c_b b^2(\log b)^k,
 \qquad 3\le k\le26,
\end{equation}
reversing lower and upper endpoints for a negative multiplier.
Outward rounding at denominator $10^{12}$ gives
Table~\ref{tab:A}.

Write $R_N(x)=\sum_{j=0}^{22}d_jx^j$. Its Bernstein representation is
\begin{equation}\label{eq:bernstein}
 \begin{split}
 R_N(x)&=\sum_{i=0}^{22}\beta_i\binom{22}{i}x^i(1-x)^{22-i},\\
 d_j&=-\sum_{k=j+4}^{26}A_k,\qquad
 \beta_i=\sum_{j=0}^{i}\frac{\binom{i}{j}}{\binom{22}{j}}d_j.
 \end{split}
\end{equation}
The conversion follows from
\[
 x^j=\sum_{i=j}^{22}\frac{\binom{i}{j}}{\binom{22}{j}}
                         \binom{22}{i}x^i(1-x)^{22-i},
\]
which is the binomial theorem after factoring out $x^j$.
All conversion factors are positive. Thus the upper bounds for $A_k$
give lower bounds for $d_j$, which in turn give lower bounds for
$\beta_i$. Table~\ref{tab:beta} lists the resulting bounds rounded down
at denominator $10^6$. Their least value is
\[
 \frac{2718387}{10^6}>\frac{271}{100}.
\]
The Bernstein basis is nonnegative and sums to one on $[0,1]$.
Consequently $R_N(x)>271/100$, and~\eqref{eq:B} proves
\begin{equation}\label{eq:final-bound}
 R(x)>\frac{271}{100}-\frac1{200}
      =\frac{541}{200}>\frac{27}{10}
 \qquad(0\le x\le1).
\end{equation}
The strict inequalities hold at the continued endpoints as well.

\begin{proof}[Completion of Theorem~\ref{thm:main}]
The $n=2$ row consists of $\Gamma_2^2=\rho_1>0$.
Lemma~\ref{lem:edges} gives both entries at $n=3$ and then the two outer
entries at $n=4$. Equation~\eqref{eq:factorizations} and $q>0$ give the
remaining entry at $n=4$. Apply Lemma~\ref{lem:edges} once more for the
outer entries at $n=5$. Its first interior entry is positive because
$N_{53}=qK>0$. The second is positive because
\eqref{eq:final-bound} gives $P(t)>0$ for $1<t<2$, and
$\Gamma_5^4=P/(16D_4)$ with $D_4>0$.
The analytic continuation and its lower bound were established above.
\end{proof}

\section{Reproducibility and limitations}

The accompanying file \texttt{verify\_certificate.py} uses only Python's
standard library. It expands the Cramer determinants with integer
polynomials, verifies the formal endpoint identities, constructs the
rational log intervals, and checks every coefficient bound used above.
It runs without external data and without writing files:
\begin{center}\texttt{python3 verify\_certificate.py}\end{center}
The finite assertions must be read together with the covariance,
zero-count and uniform remainder arguments; successful execution alone
is not a proof of those analytic steps.

A supplementary check used Arb ball arithmetic~\cite{Johansson2017}
through python-flint 0.9.0. At 192-bit precision, Taylor bounds verified
$P'>0$ on $[1,17/16]$, $P'''<0$ on $[31/16,2]$, and $P>0$ on the
112 adjacent intervals of length $1/128$ covering $[17/16,31/16]$.
Together with~\eqref{eq:endpoints}, these signs give a second interval
check of positivity. Seven direct determinant checks, including
$t=1+2^{-30}$ and $t=2-2^{-30}$, were additional finite diagnostics.
The rational certificate is the proof of the stronger bound in
Theorem~\ref{thm:main}. Rounded ball displays are not used as standalone
certificates; two stored short displays include zero even though the
recomputed underlying intervals have positive lower endpoints.

The theorem makes no assertion for $n\ge6$, for bilateral conditioning,
or about an optimal normalized lower bound. At $H=1/2$ the predictor
coefficients vanish; at $H=1$ the covariance matrix is singular.
Analytic continuation of~\eqref{eq:margin} does not define a predictor
at that singular endpoint. The outer-coefficient lemma supplies no
general propagation of the interior signs; the proof gives no
interior-sign argument for $n\ge6$.

\paragraph{Computational and AI assistance.}
OpenAI Codex assisted the proof development, literature screening,
implementation of exact and ball-arithmetic checks, and preparation of
this manuscript. A separate Codex context checked the mathematical
argument and recomputed the finite certificate using a different
symbolic determinant calculation and exponential bounds for the
logarithmic endpoints. That additional check is not external peer
review. The mathematical support consists of the arguments and exact
data given here, not of the model's assessment.

\appendix
\section{Exact certificate data}\label{app:data}

These tables specify the exponential sum and every rounded coefficient
bound used in the proof. The verifier reconstructs them from the
definitions; no floating-point values enter Tables~\ref{tab:logs}--\ref{tab:beta}.
\begin{table}[htbp]
\centering
\small
\caption{The pairs $(b,c_b)$ in $P(t)=\sum_b c_b b^t$.}\label{tab:terms}
\begin{tabular}{rr@{\qquad}rr@{\qquad}rr}
\toprule
$b$ & $c_b$ & $b$ & $c_b$ & $b$ & $c_b$ \\
\midrule
1 & -2 & 15 & 4 & 48 & 5 \\
2 & 2 & 16 & -20 & 54 & -1 \\
3 & 6 & 18 & 12 & 60 & 8 \\
4 & -11 & 24 & -12 & 64 & -6 \\
5 & -2 & 27 & -2 & 72 & -4 \\
6 & -13 & 30 & -6 & 80 & -1 \\
8 & 20 & 32 & 17 & 81 & 2 \\
9 & -4 & 36 & -1 & 90 & -1 \\
10 & 7 & 40 & -4 & 96 & 3 \\
12 & 5 & 45 & -2 & 108 & -1 \\
\bottomrule
\end{tabular}
\end{table}

\begin{table}[htbp]
\centering
\small
\caption{Integers $l_p$ such that $\log p\in[l_p,l_p+1]/10^{30}$.}\label{tab:logs}
\begin{tabular}{rr}
\toprule
$p$ & $l_p$ \\
\midrule
2 & 693147180559945309417232121458 \\
3 & 1098612288668109691395245236922 \\
5 & 1609437912434100374600759333226 \\
\bottomrule
\end{tabular}
\end{table}

\begin{table}[htbp]
\centering
\small
\caption{Integers $a_k$ such that $A_k\in[a_k,a_k+1]/10^{12}$.}\label{tab:A}
\begin{tabular}{rr@{\qquad}rr}
\toprule
$k$ & $a_k$ & $k$ & $a_k$ \\
\midrule
3 & 27182814079941 & 15 & 517016292820 \\
4 & -83550320679337 & 16 & -267302499770 \\
5 & 104183752676488 & 17 & 113465614115 \\
6 & -59505510096329 & 18 & -41923038480 \\
7 & -5716721037909 & 19 & 13880603694 \\
8 & 42148576442014 & 20 & -4191132471 \\
9 & -42718780434576 & 21 & 1167752733 \\
10 & 26787425848906 & 22 & -302833624 \\
11 & -11651227519997 & 23 & 73580734 \\
12 & 3127754437162 & 24 & -16839439 \\
13 & 74554594434 & 25 & 3645733 \\
14 & -694188830910 & 26 & -749424 \\
\bottomrule
\end{tabular}
\end{table}

\begin{table}[htbp]
\centering
\small
\caption{Integers $b_i$ giving the certified bounds $\beta_i\ge b_i/10^6$.}\label{tab:beta}
\begin{tabular}{rr@{\qquad}rr@{\qquad}rr}
\toprule
$i$ & $b_i$ & $i$ & $b_i$ & $i$ & $b_i$ \\
\midrule
0 & 27182814 & 8 & 11936093 & 16 & 5107496 \\
1 & 24620654 & 9 & 10737778 & 17 & 4594788 \\
2 & 22265492 & 10 & 9657511 & 18 & 4134485 \\
3 & 20109735 & 11 & 8684622 & 19 & 3721266 \\
4 & 18143415 & 12 & 7809147 & 20 & 3350327 \\
5 & 16355122 & 13 & 7021851 & 21 & 3017335 \\
6 & 14732704 & 14 & 6314230 & 22 & 2718387 \\
7 & 13263778 & 15 & 5678484 &  &  \\
\bottomrule
\end{tabular}
\end{table}

\clearpage
\bibliographystyle{plainurl}
\bibliography{references}
\end{document}
